\section{Two-stage stochastic formulation}\label{sec:model}

The placement is chosen before scenario $s$ is observed. Afterwards, fire
propagates along the directed arcs of $G_s^D$ and incurs the loss in
\eqref{eq:loss}. This is a two-stage problem in the decision-theoretic
sense: the second-stage burn states are determined by the first-stage
placement and the realized graph, with no additional treatment decision.
We first give a formulation for a generic scenario graph
$G_s=(V_s,E_s)$; the principal optimization experiments use $G_s^D$.

Let $x_{sk}$ indicate whether cell $k$ burns. We use the convention
$\bar y_{sj}=y_j$ if $j\in C\setminus\{r_s\}$, and
$\bar y_{sj}=0$ otherwise. A feasible propagation state satisfies
\begin{subequations}\label{eq:propagation}
\begin{align}
 x_{sr_s}&=1,\label{eq:ignite}\\
 x_{si}-x_{sj}&\le\bar y_{sj}
       &&((i,j)\in E_s),\label{eq:arcs}\\
 0\le x_{sk}&\le1 &&(k\in V_s).\label{eq:burn-domain}
\end{align}
\end{subequations}
Thus, if $i$ burns and $j$ is not a selected non-ignition firebreak,
$j$ must burn. The implemented extensive model declares $x$ binary;
the linear scenario subproblem relaxes $x$ to the interval in
\eqref{eq:burn-domain}. No explicit upper bound forcing a selected
cell's burn variable to zero is needed to obtain the correct minimum loss:
the minimum reachable-set vector has that property, as proved below.
Incoming arcs to an eligible ignition cannot change its fixed burn state,
so treating $\bar y_{sr_s}$ as zero is equivalent to the implementation's
arc expression for the purpose of the optimum.

The empirical loss of a placement in the optimization representation is
$L(y)=\sum_s\rho_s Q_s^D(y)$. To represent the upper tail, let
$\beta\in(0,1)$ be a confidence level. For finite scenario probabilities,
the Rockafellar--Uryasev representation is
\begin{equation}\label{eq:cvar}
 \operatorname{CVaR}_{\beta}(Q^D(y))
 =\min_{t\in\mathbb R}
 \left\{t+\frac{1}{1-\beta}
       \sum_{s\in\mathcal S}\rho_s[Q_s^D(y)-t]^+\right\},
 \qquad [a]^+=\max\{a,0\}.
\end{equation}
The parameter $\lambda\in[0,1]$ interpolates between the mean and this
tail measure. The risk objective is
\begin{equation}\label{eq:risk}
 \min_{y\in\mathcal Y}
  (1-\lambda)\sum_s\rho_s Q_s^D(y)
       +\lambda\operatorname{CVaR}_{\beta}(Q^D(y)).
\end{equation}
Expected loss uses $\lambda=0$, pure CVaR uses $\lambda=1$, and an
interior value gives mean--CVaR. In every case, the same first-stage
placements and scenario-loss functions are used; only the aggregation of
losses changes \citep{rockafellar2000optimization}.

For completeness, introducing nonnegative excess variables $z_s$ yields
the extensive finite-sample formulation
\begin{subequations}\label{eq:extensive}
\begin{align}
 \min_{y,x,t,z}\quad &(1-\lambda)
        \sum_s\rho_s\sum_{k\in V_s}w_kx_{sk}
       +\lambda\left(t+\frac{1}{1-\beta}
               \sum_s\rho_sz_s\right)\label{eq:ext-obj}\\
 \text{s.t.}\quad &y\in\mathcal Y,\label{eq:ext-budget}\\
 &x_{sr_s}=1 &&(s\in\mathcal S),\label{eq:ext-ignition}\\
 &x_{si}-x_{sj}\le\bar y_{sj}
       &&(s\in\mathcal S,\ (i,j)\in E_s),\label{eq:ext-graph}\\
 &z_s\ge\sum_{k\in V_s}w_kx_{sk}-t,
       \quad z_s\ge0 &&(s\in\mathcal S),\label{eq:ext-tail}\\
 &x_{sk}\in\{0,1\} &&(s\in\mathcal S,\ k\in V_s),
       \qquad t\in\mathbb R.\label{eq:ext-domains}
\end{align}
\end{subequations}
For $\lambda=0$, the $t,z$ variables and their constraints can be
omitted. The displayed formulation allows $t$ to be unrestricted, whereas
the extensive implementation bounds it between zero and a valid maximum
scenario loss. Those bounds retain an optimum because scenario losses are
nonnegative and bounded.

Fix $s$ and a binary $y\in\mathcal Y$. The scenario recourse LP minimizes
$\sum_{k\in V_s}w_kx_{sk}$ subject to
\eqref{eq:propagation}. Write its value as $q_s(y)$.
This recourse model has a particularly simple graph interpretation.

\begin{proposition}[Reachability and relatively complete recourse]
\label{prop:recourse}
For every feasible binary placement and every finite directed scenario
graph, the recourse LP is feasible and bounded. The indicator vector of
$R_s(y)$ is feasible and componentwise no larger than any other feasible
burn vector. Hence it is an optimal binary solution,
$q_s(y)=Q_s(y)$, and the model has relatively complete recourse.
\end{proposition}
\begin{proof}
The indicator of the reachable set has value one at the ignition.
If $(i,j)$ is an arc with a reachable tail and an unselected non-ignition
head, then $j$ is reachable; thus its propagation constraint holds.
All other arc constraints are immediate. Conversely, starting with
$x_{sr_s}=1$ and applying \eqref{eq:arcs} along any unblocked
ignition-to-$k$ path forces $x_{sk}=1$ for every reachable $k$.
The indicator is therefore componentwise minimal among feasible points;
nonnegative weights make it optimal. It has finite loss at most
$\sum_{k\in V_s}w_k$, establishing boundedness and recourse feasibility
for every first-stage feasible placement.
\end{proof}

Proposition~\ref{prop:recourse} justifies the scenario LP used in Benders
decomposition, even though the extensive implementation uses binary burn
variables. It also explains why one needs optimality cuts but no recourse
feasibility cuts. With zero-weight cells the LP may have other, nonminimal
optimal solutions; the reachable-set indicator still gives an optimal
solution with the intended physical interpretation.
